\section{Conditioning and absolute Newman observation}\label{sec:observation}

Load the source once, at parameter $1/2$:
\begin{equation}\label{eq:initial}
 \rho_*=|f_{1/2}\rangle\langle f_{1/2}|\otimes R,\qquad
 \lambda(t)=\tfrac12-\gamma t,\qquad 0\le t\le T.
\end{equation}
The source-position map $u=v=x_{\rm phys}/\ell_0$ and this loader remain
supplied. The prepared vectors lie in every needed finite constituent
domain by Proposition~\ref{prop:state}.

\begin{theorem}[Exact conditional family and calibrated readout]\label{thm:readout}
Let $\widetilde\rho_t=K_0(t,0)\rho_*K_0(t,0)^*$ and
$p_t=\Tr\widetilde\rho_t$. Then
\begin{equation}\label{eq:conditional}
 \widetilde\rho_t=p_tU(t)\rho^I_{\lambda(t)}U(t)^*,\qquad
 p_t=\frac{Z_{\lambda(t)}}{Z_{1/2}}\in(0,1],\qquad
 \rho_t:=\frac{\widetilde\rho_t}{p_t}=U(t)\rho^I_{\lambda(t)}U(t)^*.
\end{equation}
For real $k$ the bounded effects
$E_\pm(t,k)=[I\pm\cos(kQ_t)]/2$ give
\begin{equation}\label{eq:q}
 q_\pm(t,k)=\Tr(\widetilde\rho_tE_\pm(t,k))
 =\frac12\left[p_t\pm\frac{H_{\lambda(t)}(k)}{H_{1/2}(0)}\right].
\end{equation}
Together with $q_{\rm fail}=1-p_t$, these are a normalized probability
law, and \eqref{eq:mainintro} holds on the closed time interval.
\end{theorem}
\begin{proof}
In the interaction frame,
\[
 \ee^{-\gamma tX^2/2}f_{1/2}
 =\sqrt{Z_{\lambda(t)}/Z_{1/2}}\,f_{\lambda(t)}.
\]
This proves \eqref{eq:conditional}. Positivity and compact-slab normalization
give $p_t\ge Z_{-1/2}/Z_{1/2}>0$; the attenuation is a contraction, so
$p_t\le1$. Real spectral calculus gives $0\le E_\pm\le I$ and
$E_++E_-=I$. Bounded-factor trace cyclicity and the complete angular
trace $\Tr R=1$ yield
\begin{equation}\label{eq:cosine}
 C(t,k):=\Tr(\rho_t\cos(kQ_t))
 =\int_\R\frac{F_{\lambda(t)}(v)}{Z_{\lambda(t)}}\cos(kv)\dd v
 =\frac{H_{\lambda(t)}(k)}{H_{\lambda(t)}(0)}.
\end{equation}
Multiplying by $p_t$ proves \eqref{eq:q}. On the initial state, the three
effects are $K_0^*E_+K_0$, $K_0^*E_-K_0$ and $I-K_0^*K_0$; they sum to
$I$ and are positive. Thus the quantities are actual joint probabilities
of the stipulated instrument and no-count readout, with the standard
operation normalization \cite[Section III, (17)--(19)]{JacobsSteck2006v1}.
\end{proof}

The success frequency is essential to absolute calibration. Selected
successful samples give $H_\lambda(k)/H_\lambda(0)$; multiplying by $p_t$
and the single initial scale gives $H_\lambda(k)$. At $k=0$,
$q_+=p_t$ and $q_-=0$. At $t=1/(2\gamma)$ the readout is precisely
\eqref{eq:xi}. These checks retain the full-line factor $Z_\lambda=2H_\lambda(0)$.
The observed density is the distribution of $Q_t$, not generally the
laboratory position $X$ after free dispersion.

\subsection{Regular state motion and signed evolution}
Write $w_t=|f_{\lambda(t)}|^2$ and $\mu_r(t)=\int v^r w_t(v)\dd v$.
Differentiation of the paid source integrals gives
\begin{align}
 \partial_tw_t&=-\gamma(v^2-\mu_2)w_t,
 &\dot p_t&=-\gamma\mu_2p_t,\label{eq:selection}\\
 \partial_tf_{\lambda(t)}&=-\frac\gamma2(X^2-\mu_2)f_{\lambda(t)},
 &\dot\mu_2&=-\gamma(\mu_4-\mu_2^2).\label{eq:tangent}
\end{align}
The first equation in \eqref{eq:tangent} includes both the square-root
half factor and the normalization subtraction. In particular
\begin{equation}\label{eq:speed}
 \norm{\partial_tf_{\lambda(t)}}_2^2
 =\frac{\gamma^2}{4}(\mu_4-\mu_2^2)<\infty.
\end{equation}
All fixed moments are bounded on the slab. The complete prepared density
operator is smooth to every fixed finite order in trace norm: differentiate
the rank-one ensemble and use the corresponding radial norms and summable
angular graph moments. This does not assert norm differentiability of
$U(t)$ on arbitrary states.

For a fixed polynomially bounded test function used here, the same
calculation gives
\begin{equation}\label{eq:covariance}
 \frac{d}{dt}\E_t a(X)=-\gamma\Cov_t(a(X),X^2).
\end{equation}
Consequently, with $d_0(t,k)=p_tC(t,k)$,
\begin{equation}\label{eq:heat}
 C_t=\gamma(C_{kk}+\mu_2C),\qquad
 (d_0)_t=\gamma(d_0)_{kk},\qquad
 d_0=\frac{H_{\lambda(t)}(k)}{H_{1/2}(0)}.
\end{equation}
The normalization term is absent only in the absolute contrast $d_0$.
The physical-time heat sign is positive because $\lambda(t)$ decreases.
The Fourier amplitude of $f_\lambda$, its momentum probability, and the
Fourier transform of $|f_\lambda|^2$ are different objects; only the last
produces this density-cosine observation. Entire continuation of
\eqref{eq:Hmixed} supplies no bounded nonreal-$k$ measurement.
